% ============================================================
% APPENDICE F — Progetto Finale Autonomo: BookShelf
% ============================================================
\chapter{Progetto Finale Autonomo: BookShelf}

\section*{Il Tuo Esame di Architettura}

Questo progetto non ha una guida passo-passo. Hai il brief, il \texttt{\_CONTEXT.md} di partenza, e il \textbf{workflow professionale} della Sezione~3.8 --- lo stesso metodo di decomposizione che useresti in un progetto reale. Il tuo compito è applicarlo dall'inizio alla fine: generare le specifiche, derivare i casi d'uso, strutturare le epic, scomporre in task e implementare.

Se riesci a completare BookShelf da solo, hai interiorizzato il metodo ADLC.

% ---------------------------------------------------------
\section{Il Brief}

\textbf{BookShelf} è un'applicazione per gestire la tua libreria personale. Puoi catalogare i libri che possiedi, quelli che stai leggendo e quelli che vuoi leggere. L'applicazione ha un backend API, un frontend web e un'app mobile.

\subsection{Requisiti funzionali}

\begin{enumerate}
  \item \textbf{Catalogo libri}: CRUD completo (titolo, autore, anno, genere, ISBN opzionale, copertina URL opzionale)
  \item \textbf{Stati di lettura}: Ogni libro ha uno stato --- ``Da leggere'', ``In lettura'', ``Letto'', ``Abbandonato''
  \item \textbf{Note personali}: L'utente può aggiungere note e una valutazione (1--5 stelle) a ogni libro
  \item \textbf{Ricerca e filtri}: Ricerca per titolo/autore, filtro per stato e genere
  \item \textbf{Statistiche}: Dashboard con numero libri per stato, libri letti per anno, genere più letto
  \item \textbf{Autenticazione}: Login OAuth~2.0 (Google)
  \item \textbf{Multi-utente}: Ogni utente vede solo la propria libreria
\end{enumerate}

\subsection{Requisiti non funzionali}

\begin{itemize}
  \item Backend: Node.js + Express + PostgreSQL + Prisma
  \item Frontend: React + Tailwind CSS
  \item Mobile: Flutter
  \item Deploy: piattaforme cloud a scelta (Railway, Vercel, Neon o equivalenti)
\end{itemize}

% ---------------------------------------------------------
\section{Il \texttt{\_CONTEXT.md} di Partenza}

Questo è il tuo punto di partenza. Copialo nella root del progetto e usalo come base. Puoi --- e dovresti --- modificarlo e migliorarlo man mano che procedi.

\begin{lstlisting}[language={},title={\_CONTEXT.md --- BookShelf}]
# Progetto: BookShelf

## Scopo
Applicazione web + mobile per la gestione di una libreria
personale. Ogni utente puo catalogare libri, tracciare lo stato
di lettura, aggiungere note e visualizzare statistiche.

## Tecnologie
- Backend: Node.js 20+, Express.js, PostgreSQL 16, Prisma ORM
- Frontend: React 18+ con Vite, Tailwind CSS
- Mobile: Flutter 3.x con Dart, Riverpod
- Auth: OAuth 2.0 con Google (passport.js)
- Token: JWT (access 1h, refresh 7d)

## Struttura
bookshelf/
+-- _CONTEXT.md
+-- backend/
|   +-- src/
|   |   +-- routes/
|   |   +-- controllers/
|   |   +-- services/
|   |   +-- middleware/
|   |   +-- utils/
|   +-- prisma/
|   |   +-- schema.prisma
|   +-- tests/
+-- frontend/
|   +-- src/
|   |   +-- components/
|   |   +-- pages/
|   |   +-- hooks/
|   |   +-- services/
|   |   +-- context/
|   +-- tests/
+-- mobile/
    +-- bookshelf_app/
        +-- lib/
            +-- models/
            +-- providers/
            +-- screens/
            +-- services/
            +-- widgets/

## Modello Dati (schema Prisma indicativo)
- User: id, email, name, avatar, createdAt
- Book: id, userId, title, author, year, genre,
        isbn?, coverUrl?,
        status (TO_READ | READING | READ | ABANDONED),
        rating (1-5)?, notes?, createdAt, updatedAt
- ReadingSession: id, bookId, startDate, endDate?,
        pagesRead?

## Convenzioni
- Naming: camelCase per JS/TS, snake_case per DB columns
- Risposte API: { success: boolean, data: T, error?: string }
- Validazione: Zod su tutti gli endpoint
- Commit: Conventional Commits

## Vincoli (SEC/PERF)
- SEC-01: JWT secret >= 256 bit, in variabile d'ambiente
- SEC-02: Ogni query filtra per userId -- NESSUNA eccezione
- SEC-03: Input sanitizzato e validato prima del database
- SEC-04: Rate limiting su endpoint auth (10 req/min)
- SEC-05: Helmet.js per header di sicurezza
- PERF-01: Paginazione su GET /books (default 20, max 100)
- PERF-02: Indice su Book.userId e Book.status

## Comandi
- Backend dev: cd backend && npm run dev
- Backend test: cd backend && npm test
- Frontend dev: cd frontend && npm run dev
- Frontend test: cd frontend && npm test
- DB migrate: cd backend && npx prisma migrate dev
- Mobile run: cd mobile/bookshelf_app && flutter run
- Mobile test: cd mobile/bookshelf_app && flutter test
\end{lstlisting}

% ---------------------------------------------------------
\section{Il Workflow: Dall'Idea al Codice}

Questo progetto non ha una guida passo-passo, ma hai un \textbf{metodo}. Applica il workflow professionale della Sezione~3.8 per scomporre BookShelf in unità di lavoro gestibili --- esattamente come faresti in un progetto reale.

\subsection{Passo 1: Genera le specifiche}

Apri una chat con l'IA e usa il brief come input:

\begin{lstlisting}[language={}]
"Ecco il brief di BookShelf: [incolla F.1].
 Genera un PRD completo con requisiti funzionali (RF-01, RF-02...),
 requisiti non funzionali (RNF-01...), vincoli tecnici e priorita'.
 Salva in docs/PRD.md."
\end{lstlisting}

\subsection{Passo 2: Deriva i casi d'uso}

\begin{lstlisting}[language={}]
"A partire dal PRD in docs/PRD.md, genera i casi d'uso principali.
 Per ogni caso d'uso: attore, precondizione, flusso principale,
 flussi alternativi, postcondizione. Salva in docs/USE_CASES.md."
\end{lstlisting}

\subsection{Passo 3: Progettazione e scelte tecnologiche}

Prima di scomporre in epic e task, fai scegliere all'umano (cioè a te) le decisioni architetturali:

\begin{lstlisting}[language={}]
"A partire dal PRD e dai casi d'uso di BookShelf, proponi l'architettura
 del sistema. Per ogni decisione tecnica (database, ORM, autenticazione,
 struttura progetto, pattern API) presenta 2-3 opzioni con pro, contro
 e raccomandazione. Io scelgo, tu documenti in docs/DESIGN.md."
\end{lstlisting}

Prendi ogni decisione consapevolmente --- l'IA propone, tu decidi. Il risultato è un \texttt{docs/DESIGN.md} con le ADR (Architecture Decision Records) che guideranno tutto il progetto.

\subsection{Passo 4: Struttura le epic}

\begin{lstlisting}[language={}]
"A partire dai casi d'uso, organizza il lavoro in epic (E-001, E-002...).
 Per ogni epic: titolo, obiettivo, casi d'uso coperti, dipendenze.
 Salva l'indice in docs/epics/INDEX.md e un file per epic."
\end{lstlisting}

Le epic suggerite coprono le 4~fasi del progetto:

\begin{center}
\begin{tabularx}{\textwidth}{l l X}
\toprule
\textbf{Fase} & \textbf{Epic} & \textbf{Capitoli} \\
\midrule
Backend & E-001 Setup, E-002 Auth, E-003 CRUD, E-004 Filtri, E-005 Stats & Cap.~6--8 \\
Frontend & E-006 Frontend Web & Cap.~9--10 \\
Mobile & E-007 App Mobile & Cap.~11--12 \\
Produzione & E-008 Testing/Sicurezza, E-009 Deploy & Cap.~14--15 \\
\bottomrule
\end{tabularx}
\end{center}

\subsection{Passo 5: Scomponi in task}

Per ogni epic, chiedi all'IA di generare task atomici:

\begin{lstlisting}[language={}]
"Scomponi l'epic E-003 in task atomici (T-003.1, T-003.2...).
 Ogni task deve avere: acceptance criteria verificabili,
 file coinvolti, vincoli SEC/PERF applicabili, story point."
\end{lstlisting}

Segui il template task della Sezione~3.8. Salva in \texttt{docs/epics/tasks/E-003/}.

\subsection{Passo 6: Implementa task per task}

\begin{lstlisting}[language={}]
"Implementa T-003.1. Leggi docs/epics/tasks/E-003/T-003.1_*.md.
 Rispetta i vincoli dal _CONTEXT.md."
\end{lstlisting}

Ogni sessione chiude uno o più task. Al termine, aggiorna lo stato e il \texttt{\_CONTEXT.md}.

\begin{tipbox}
Non procedere alla Fase~2 (Frontend) finché tutti i task backend non sono completati e i test non passano. La decomposizione in task ti dà visibilità --- usala.
\end{tipbox}

% ---------------------------------------------------------
\section{Le Sfide per Fase}

Ecco gli obiettivi di alto livello per ogni fase. La scomposizione dettagliata in epic e task è compito tuo --- usando il workflow del Passo~4 e~5.

\subsection{Fase 1: Backend (Capitoli 6--8)}

\begin{enumerate}
  \item Inizializza il progetto Node.js con Express
  \item Configura Prisma con lo schema \texttt{Book} e \texttt{User}
  \item Implementa gli endpoint CRUD per i libri
  \item Aggiungi autenticazione OAuth~2.0 con Google
  \item Aggiungi l'endpoint statistiche (\texttt{GET /api/stats})
  \item Scrivi test unitari e di integrazione
  \item Verifica la security checklist (SEC-01 $\rightarrow$ SEC-05)
\end{enumerate}

\subsection{Fase 2: Frontend (Capitoli 9--10)}

\begin{enumerate}
  \item Crea il progetto React con Vite e Tailwind
  \item Implementa la pagina login con redirect OAuth
  \item Crea la dashboard con lista libri e filtri
  \item Implementa il form per aggiungere/modificare libri
  \item Crea la pagina statistiche con conteggi e grafici semplici
  \item Integra con il backend reale (rimuovi mock)
  \item Scrivi test con testing-library
\end{enumerate}

\subsection{Fase 3: Mobile (Capitoli 11--12)}

\begin{enumerate}
  \item Crea il progetto Flutter
  \item Implementa l'autenticazione da mobile
  \item Crea la lista libri con pull-to-refresh
  \item Implementa il dettaglio libro con note e rating
  \item Aggiungi ricerca e filtri
  \item Scrivi test widget
\end{enumerate}

\subsection{Fase 4: Produzione (Capitoli 14--15)}

\begin{enumerate}
  \item Esegui audit OWASP Top~10 sul backend
  \item Applica i fix di sicurezza
  \item Deploy backend su Railway
  \item Deploy frontend su Vercel
  \item Database su Neon
  \item Verifica che tutto funzioni end-to-end
\end{enumerate}

% ---------------------------------------------------------
\section{Criteri di Successo}

Il progetto è completato quando:

\begin{center}
\begin{tabularx}{\textwidth}{c X}
\toprule
\textbf{Stato} & \textbf{Criterio} \\
\midrule
$\square$ & Un utente può fare login con Google \\
$\square$ & Può aggiungere, modificare e cancellare libri \\
$\square$ & Può cambiare lo stato di lettura \\
$\square$ & Può aggiungere note e rating \\
$\square$ & Può cercare e filtrare i libri \\
$\square$ & Vede le statistiche sulla dashboard \\
$\square$ & L'app mobile mostra gli stessi dati \\
$\square$ & Tutti i test passano \\
$\square$ & La security checklist è completata \\
$\square$ & L'applicazione è accessibile online (deployata) \\
$\square$ & Il \texttt{\_CONTEXT.md} è aggiornato con le lezioni apprese \\
\bottomrule
\end{tabularx}
\end{center}

% ---------------------------------------------------------
\section{Suggerimenti}

\begin{itemize}
  \item \textbf{Segui il workflow della Sezione~3.8.} Prima genera PRD, poi use case, poi epic, poi task. Solo dopo implementa.
  \item \textbf{Inizia dal backend.} È la base di tutto il resto.
  \item \textbf{Usa il \texttt{\_CONTEXT.md} fin dall'inizio.} Non procedere senza contesto.
  \item \textbf{Testa ogni fase} prima di passare alla successiva. Non costruire il frontend su un backend che non funziona.
  \item \textbf{Se l'IA genera qualcosa che non capisci}, chiedi: ``Spiega questo codice riga per riga'' --- è sempre un prompt valido.
  \item \textbf{Se ti blocchi}, riguarda il capitolo corrispondente del libro. La tecnica è la stessa --- cambia solo il dominio applicativo.
  \item \textbf{Aggiorna il \texttt{\_CONTEXT.md}} ogni volta che scopri qualcosa di nuovo. Questo è l'Apprendimento Continuo (Fase~6 dell'ADLC) in azione.
\end{itemize}

\vfill

\begin{center}
\textit{Se hai completato BookShelf, complimenti: sei un architetto di contesto. Il metodo ADLC non è più qualcosa che hai letto --- è qualcosa che sai fare.}
\end{center}
